\documentclass[conference]{IEEEtran}

\usepackage{cite}
\usepackage{amsmath,amssymb,amsfonts}
\usepackage{graphicx}
\usepackage{textcomp}
\usepackage{xcolor}
\usepackage{booktabs}
\usepackage{url}
\usepackage{array}
\usepackage{stfloats}
\usepackage{float}

\begin{document}

\title{
\centering
\fontsize{24}{28}\selectfont
\textbf{P4-Based Intrusion Detection System with Rule}\\[0.15cm]
\textbf{Distillation for Data Plane Security}
}

\author{
\IEEEauthorblockN{
\fontsize{12}{14}\selectfont
Venugopal Rao Ponnamaneni \hspace{2.5cm} Moditha Lekkala
}
\vspace{0.15cm}
\IEEEauthorblockA{
\fontsize{10}{12}\selectfont
Saint Louis University
}
}

\maketitle

\begin{abstract}
Modern network infrastructures require intrusion
detection systems that can identify malicious traffic in real time
without adding heavy latency or host-side processing overhead.
Traditional intrusion detection systems such as software-based
monitors inspect traffic outside the forwarding path, which
can reduce throughput and increase CPU load. This project
implements a P4-based intrusion detection system that performs
classification directly inside the programmable data plane. The
system uses machine-learning-derived rule distillation, where
decision-tree behavior from CICIDS2017 traffic is converted into
compact ternary match-action rules and deployed on a BMv2
switch. The switch aggregates per-flow statistics over a two-
second window and decides whether to forward or drop traffic.
The final system achieved 86\% accuracy on 100 test flows,
with 100\% precision, 72\% recall, and zero false positives. The
results show that learned traffic behavior can be transformed into
efficient switch rules for low-latency intrusion detection while
preserving legitimate traffic.
\end{abstract}

\section{INTRODUCTION}

Modern networks face a growing number of security threats,
including scanning, denial-of-service behavior, and malicious
traffic patterns that must be detected quickly. Traditional
intrusion detection systems usually run in software on a host
or monitoring server. Although these systems can be effective,
they often require packets to leave the forwarding path for
inspection. This increases latency, consumes CPU resources,
and may limit throughput when traffic volume becomes high.

Programmable data planes provide a different approach.
Using the P4 language, packet-processing behavior can be
programmed directly into a switch pipeline. This makes it
possible to perform lightweight security decisions inside the
network device instead of forwarding all suspicious traffic to
an external analyzer. The purpose of this project is to test
whether machine-learning-derived intrusion detection logic
can be represented as compact rules that execute efficiently
inside a programmable switch.

This project implements a real-time intrusion detection
system using P4, BMv2, and rule distillation. The system
uses CICIDS2017-style benign and malicious traffic profiles,
extracts flow-window features, and classifies traffic using
ternary match-action rules. The project addresses the gap
between traditional software-based intrusion detection and
the need for low-latency security enforcement directly inside
programmable network devices. The main goal is to preserve
benign traffic while dropping malicious flows with minimal
latency while maintaining reproducible and hardware-friendly
classification behavior.

A key challenge in programmable data-plane security is
the mismatch between the expressiveness required by machine
learning models and the execution constraints of switch
pipelines. Switches operate under strict latency budgets and
support only a limited instruction set. This project bridges that
gap by translating offline-learned behavior into a small,
hardware-compatible rule set, demonstrating that meaningful
security logic can survive the translation from a high-level
model representation to a deployable data-plane policy.

\section{RELATED WORK}

Intrusion detection systems have traditionally used
signature-based methods or machine learning models
deployed in software. Tools such as Snort and Suricata
inspect packets outside the switch pipeline. These systems are
useful for deep inspection, but they can introduce processing
overhead when traffic rates are high.

Recent work has explored programmable switches for
security. P4ID showed that intrusion detection logic can be
moved into the programmable data plane, reducing the amount
of traffic that must be processed by external IDS systems.
This project follows the same general direction but focuses on
converting machine-learning-derived rules into a small set of
ternary match-action rules that can run inside BMv2.

Machine learning models such as decision trees are useful
for classification because their decision paths can be converted
into interpretable rules. However, directly deploying complex
models in the data plane is difficult because programmable
switches have limited computation and memory. Rule distilla-
tion addresses this limitation by converting learned behavior
into simple match-action entries. Work such as Quark has
demonstrated that even convolutional neural networks can be
partially mapped to programmable pipelines, though such
approaches require careful approximation to fit within switch
resource budgets.

The CICIDS2017 dataset has been widely used as a
benchmark for evaluating network intrusion detection systems.
It provides labeled packet captures of both benign and
malicious traffic, including port scanning, DDoS, and botnet
activities. Using flow-level features derived from this dataset
provides a realistic foundation for training and validating
classification rules in a controlled environment.

\section{PROBLEM DEFINITION AND SOLUTION DESIGN}

The main problem is to detect malicious network traffic
in real time without relying completely on a host-side IDS.
The system must classify traffic directly inside the switch
pipeline while preserving high throughput and low latency.
A practical data-plane IDS should avoid dropping legitimate
traffic, because false positives can interrupt normal network
operation.

The proposed solution has four stages. First, traffic behavior
is learned offline from CICIDS2017-style data using decision-
tree-based logic. Second, the learned behavior is converted
into five ternary match-action rules stored in \texttt{rules.json}.
Third, the P4 switch aggregates flow-window statistics for
incoming traffic over a two-second window. Finally, the switch
applies the decision table and either forwards benign traffic or
drops malicious traffic.

The P4 switch tracks five main features for each flow:
maximum packet length, minimum packet length, TCP PSH
flag count, total packet length, and packet count. Benign
traffic is represented by larger packet sizes around 500 bytes
and mixed PSH behavior. Malicious traffic is represented by
smaller packet sizes around 60 bytes and no PSH flags. These
features are sufficient for the controlled benchmark used in
this project.

The system architecture separates the control plane from
the data plane. The control plane, implemented in Python,
is responsible for loading classification rules into the switch
table at startup using \verb|simple_switch_CLI|. Once rules are
installed, the data plane handles all per-packet and per-flow
decisions autonomously without requiring further control-plane
involvement. This separation ensures that the classification
latency measured during benchmarking reflects only the
data-plane cost and not any runtime coordination overhead.
The two-second aggregation window is chosen to balance
detection responsiveness against feature stability, as very
short windows can lead to incomplete flow statistics while
very long windows delay the mitigation action.

\section{PROJECT OBJECTIVE}

The purpose of this project is to demonstrate that intrusion
detection logic derived from machine learning models can be
efficiently deployed inside programmable network switches.
The project is intended for researchers, networking students,
and practitioners interested in programmable data planes,
low-latency network security, and deployable AI-assisted
network defense mechanisms.

The project provides a reproducible implementation using
P4, BMv2, Docker, Python automation scripts, and rule-based
classification derived from offline learning. The system is
designed to show how programmable switches can support
practical real-time intrusion detection without relying entirely
on external software-based monitoring systems.

Specifically, the project targets three measurable goals:
achieving greater than 80\% classification accuracy on a
controlled 100-flow benchmark, maintaining zero false posi-
tives so that legitimate traffic is never dropped, and executing
per-flow classification in under 100 microseconds to demon-
strate the latency advantage of data-plane enforcement over
software-based IDS alternatives.

\section{METHODOLOGY}

The project uses a practical form of knowledge distillation.
Instead of executing a full machine learning model online,
the offline learned behavior is converted into a compact set of
switch-compatible rules. This allows the data plane to perform
classification with simple table lookups instead of expensive
model inference.

\subsection{Knowledge Distillation}

The project adopts a practical form of knowledge distilla-
tion, where a richer offline machine learning representation is
transformed into a compact and deployable rule-based model
suitable for execution in the programmable data plane. Instead
of executing a full decision tree or complex model at runtime,
the learned decision boundaries are converted into a small set
of match-action rules.

This enables the switch to perform classification directly at
line rate using simple table lookups. As a result, the system
avoids the overhead of executing a full machine learning model
during packet processing while still preserving meaningful
classification accuracy. This approach is essential for pro-
grammable switches where latency and resource constraints
are strict.

The distillation process begins by training a decision tree
classifier on CICIDS2017-style flow records. The tree's leaf
paths are then analyzed to extract the feature thresholds that
most effectively separate benign from malicious classes. These
thresholds are encoded as ternary value-mask pairs, where the
mask controls which bits of the feature register participate in
the match. A small set of five rules is sufficient to represent
the key decision boundaries for the traffic profiles used in this
benchmark. The resulting \texttt{rules.json} file is loaded into
the P4 table at startup, and no further model interaction is
required during operation.

\subsection{TCAM and Ternary Matching}

The deployed rules use ternary matching, where each rule
consists of a value and a mask, corresponding to TCAM-
style behavior used in network hardware. Exact matches are
achieved using all-one masks, while wildcard behavior is
enabled by masking specific bits. Partial masks allow range-
like grouping of values, enabling multiple decision boundaries
to be represented compactly.

This approach allows complex classification logic to be
compressed into a small number of rules. Rule priorities are
used to resolve overlaps between general and specific con-
ditions, ensuring correct classification decisions. As a result,
the system produces a compact and hardware-efficient policy
derived from learned behavior.

In practice, a ternary rule for detecting malicious short-
packet flows encodes a value such as 60 for the maximum
packet length feature with a mask that matches the upper
portion of the value range. When incoming flow statistics
are compared against this entry, a match triggers the drop
action. Rules with higher priority are evaluated first, so
specific conditions such as an exact packet-length match take
precedence over broader wildcard entries. The use of TCAM-
style semantics in BMv2 allows this priority resolution to
occur at lookup time without additional branching logic in
the P4 program itself.

\subsection{Flow Feature Aggregation}

Flow-level features are maintained using P4 registers indexed
by a hash of the five-tuple: source IP address, destination IP
address, source port, destination port, and protocol. Each
register array stores a running aggregate for one feature,
including the maximum observed packet length, minimum
observed packet length, total packet length, packet count,
and PSH flag count. At each packet arrival, the ingress
pipeline reads the current register values, updates them with
the new packet's fields, and writes the updated values back.
After the two-second observation window expires, the control
plane reads the register values using a Python polling script
and presents them as a feature vector to the classification
table. If the table matches a malicious rule, subsequent
packets belonging to that flow are dropped.

\section{CENTROID-BASED CLASSIFICATION MODEL}

\subsection{Model Overview}

In addition to rule distillation, the project incorporates a
centroid-based binary classification approach optimized for the
P4 data plane. This model provides a lightweight alternative
to traditional machine learning models and is well-suited for
execution in programmable switches.

\subsection{Intuition and Logic}

The centroid-based model represents each class, namely
normal and attack traffic, using a prototype bit string called
a centroid. These centroids capture the key characteristics of
each class in a simplified representation. A centroid is
computed offline by taking the bitwise majority of all training
samples belonging to a given class across each bit position.
This produces a compact binary template that approximates
the average behavior of each traffic class.

During inference, the switch compares incoming flow features
against the stored centroids using bitwise XNOR operations
to identify matching bits. A population count operation is
then applied to compute a similarity score by counting the
number of matching bits.

The final classification decision is made by selecting the
class with the higher similarity score. This method allows
efficient classification using simple bitwise operations, making
it highly suitable for real-time data-plane execution without
requiring complex computations. Because XNOR and popcount
operations map naturally to register reads and arithmetic
actions available in P4, the centroid model avoids the need
for any floating-point computation or multi-step inference
pipeline. The trade-off is that the centroid representation loses
information about feature distribution variance, which can
reduce recall when attack traffic deviates from the learned
prototype. However, the model's simplicity makes it an
attractive fallback when the ternary rule table does not
produce a confident match.

\section{IMPLEMENTATION DETAILS}

The system is implemented using P4 and BMv2. The
main P4 program is located in \texttt{src/p4/main.p4}, and the
compiled BMv2 output is stored as \texttt{main.json}. The switch
maintains registers for per-flow statistics and uses a hash range
of 8192 entries for flow tracking.

The core Python scripts support testing and automation. The
\texttt{inject.py} script generates 100 test flows in two phases,
including 50 benign flows and 50 malicious flows. The
\texttt{verify.py} script captures egress traffic for 15 seconds
and builds the confusion matrix. The \texttt{trigger.py} script
loads \texttt{rules.json} into the P4 table using
\verb|simple_switch_CLI|. The \verb|sweep_len.py| script is
used for threshold calibration.

The project repository is organized into directories for
source code, P4 programs, Python scripts, data, packet
captures, plots, Docker files, documentation, and notebooks.
The benchmark can be executed using the provided PowerShell
automation script \verb|run_full_benchmark.ps1|. Benign and
malicious flows were injected in controlled phases using
Python-based traffic generation, while the BMv2 switch
performed real-time classification using deployed ternary
match-action rules. The generated plots and metrics can be
reproduced using the provided Docker environment and
benchmark automation scripts.

The Docker environment pins all package versions to
ensure reproducibility across different host systems. The
container image includes Python 3.11, Scapy for packet
generation, and the BMv2 simple\_switch binary compiled
from the official P4 behavioral-model repository. Network
namespaces inside the container simulate the ingress and
egress interfaces of the switch. The control-plane scripts
communicate with the switch using the \verb|simple_switch_CLI|
gRPC interface to install rules and read register values. This
containerized setup eliminates dependency on any specific
hardware or operating system configuration and allows the
benchmark to be reproduced by running a single automation
script.

\section{EXPERIMENTAL EVALUATION}

The evaluation uses 100 test flows, consisting of 50 malicious
flows and 50 benign flows. The malicious flows are expected
to be dropped, while the benign flows are expected to be
forwarded. The system correctly dropped 36 malicious flows
and correctly forwarded 50 benign flows. It incorrectly
forwarded 14 malicious flows and dropped zero benign flows.

\begin{table}[H]
\caption{Confusion Matrix}
\label{tab:confusion_matrix}
\centering
\begin{tabular}{lcc}
\toprule
 & \textbf{Predicted DROP} & \textbf{Predicted FORWARD} \\
\midrule
\textbf{Actual DROP} & 36 & 14 \\
\textbf{Actual FORWARD} & 0 & 50 \\
\bottomrule
\end{tabular}
\end{table}

As shown in Table~\ref{tab:confusion_matrix}, the system produced zero false
positives. This is important because it means no legitimate
traffic was incorrectly dropped during the test.

The core evaluation metrics are computed as follows.
Accuracy is 86/100, or 86\%. Precision is 1.000 because
every predicted drop was actually malicious. Recall is 0.720
because the system detected 36 out of 50 malicious flows.
The F1 score, which balances precision and recall, is
approximately 0.837, reflecting the trade-off between the
system's conservative rule set and its strong avoidance of
false positives.

\begin{table}[H]
\caption{Summary of Evaluation Metrics}
\label{tab:metrics}
\centering
\begin{tabular}{lc}
\toprule
\textbf{Metric} & \textbf{Value} \\
\midrule
Accuracy  & 86.0\% \\
Precision & 1.000 \\
Recall    & 0.720 \\
F1 Score  & 0.837 \\
False Positives & 0 \\
\bottomrule
\end{tabular}
\end{table}

Fig.~\ref{fig:roc} shows the ROC behavior of the system. The selected
operating point prioritizes zero false positives while still
detecting most malicious traffic.

Fig.~\ref{fig:latency} shows that the P4-based approach provides much
lower latency than software IDS processing. The benchmark
reports approximately 12 microseconds per flow for the
P4-based system compared with approximately 850 microseconds
for software IDS processing, representing a speedup of
roughly 70 times. This difference arises because the P4
data plane performs classification entirely through table
lookups in the forwarding pipeline, whereas software IDS
tools must copy packets to userspace, apply pattern matching
across the full payload, and return a verdict through the
operating system network stack.

Fig.~\ref{fig:timeline} shows that detection delay is mainly caused by
the feature aggregation window and benchmark timing. The
benchmark reports a mitigation delay of approximately 4.8
seconds, which includes the windowing and verification
process rather than expensive runtime inference. The dominant
component of this delay is the two-second observation window
required to accumulate stable flow statistics, followed by
the 15-second egress capture window used by \texttt{verify.py}.
In a production deployment, the observation window could be
reduced to sub-second intervals at the cost of reduced feature
stability, and the capture window could be eliminated entirely
by triggering verification on a per-flow basis.

\begin{figure}[!t]
\centering
\includegraphics[width=0.45\textwidth]{Picture1.png}
\caption{ROC curve showing the relationship between true positive rate and false positive rate.}
\label{fig:roc}
\end{figure}

\begin{figure}[!t]
\centering
\includegraphics[width=0.45\textwidth]{Picture2.png}
\caption{Latency and throughput comparison between the P4-based IDS and software IDS processing.}
\label{fig:latency}
\end{figure}

\begin{figure}[!t]
\centering
\includegraphics[width=0.45\textwidth]{Picture3.png}
\caption{Mitigation timeline showing detection delay from traffic observation to final action.}
\label{fig:timeline}
\end{figure}

\section{DISCUSSION}

The results show that rule distillation is a practical method
for deploying intrusion detection logic in a P4 data plane.
The strongest result is zero false positives, meaning the
system did not incorrectly drop legitimate traffic. This is
important because in real deployments, false positives can
interrupt normal users and services and may trigger costly
manual review processes or rollback procedures.

The main limitation is recall. The system detected 36 out
of 50 malicious flows, leaving 14 malicious flows forwarded.
This suggests that the current rule set is conservative. The
14 missed flows likely correspond to attack traffic that does
not closely match the packet-length and PSH-flag patterns
captured by the five deployed rules. Two primary causes
can account for this gap: first, the distillation process
may have discarded decision boundaries that were useful in
the original decision tree but could not be expressed as
ternary rules without a higher false-positive risk; second,
the malicious traffic generator may have produced a subset
of flows with packet sizes or timing patterns that overlap
with benign profiles. Future work should improve recall by
adding more malicious signatures, tuning threshold masks,
expanding the feature set, and testing on a wider range of
CICIDS2017 traffic classes.

Another limitation is that the benchmark uses controlled
traffic profiles. Real networks contain more diverse packet
sizes, timing patterns, and protocol behaviors. A stronger
evaluation would include more attack categories, background
traffic, and longer test windows. However, the current project
demonstrates a working end-to-end pipeline from traffic
generation to P4 data-plane classification and metric reporting.

From an ethical perspective, the system is designed to
improve network security and reduce response time. It does
not inspect private payload content in the current design; it
uses flow-level packet statistics. Still, any real deployment
should document what features are collected, how logs are
stored, and how false positives are handled. Network-level
enforcement mechanisms that automatically drop traffic carry
the risk of being abused as a denial-of-service tool if rule
updates are not properly authenticated and audited.

\section{FUTURE WORK}

Future work will focus on improving recall while preserving
the current zero false-positive behavior. Additional malicious
traffic categories from CICIDS2017 and other benchmark
datasets can be incorporated to improve classification coverage.
One concrete direction is to expand the feature set to include
inter-arrival time statistics and byte-entropy estimates, which
can help distinguish encrypted malicious traffic from
legitimate encrypted sessions.

Future versions of the system may also support adaptive rule
updates, dynamic threshold tuning, and integration with
hardware programmable switches beyond BMv2 simulation.
Adaptive updates would allow the control plane to refresh
the match-action table at runtime based on observed traffic
distributions, enabling the system to respond to concept drift
without requiring a full offline retraining cycle.

Another important direction is extending the system toward
production-scale deployment using larger traffic traces,
background network activity, and multi-stage detection
pipelines. A multi-stage approach could combine a
lightweight first-pass ternary rule table for high-confidence
classifications with a more expressive secondary stage,
implemented in software, for ambiguous flows. This hybrid
architecture would preserve the latency advantage of
data-plane classification for the majority of traffic while
retaining the flexibility of software analysis for edge cases.
The project may also be expanded into a publishable research
platform by evaluating additional machine learning models,
rule compression strategies, and hardware-aware optimization
techniques targeting P4-capable ASICs such as the Intel
Tofino.

\section{CONCLUSION}

This project demonstrates a working P4-based intrusion
detection system that performs rule-based machine learning
classification directly in the programmable data plane. The
system converts learned behavior into ternary match-action
rules, tracks flow-window features, and executes forward or
drop decisions inside BMv2. The final benchmark achieved
86\% accuracy, 100\% precision, 72\% recall, and zero false
positives. These results show that data-plane intrusion
detection is a promising direction for low-latency network
security. The 70x latency reduction over software IDS
processing and the zero false-positive rate together
demonstrate that meaningful security enforcement can be
achieved inside a programmable switch without sacrificing
legitimate traffic. Future improvements should focus on
increasing recall, testing more attack types, and supporting
adaptive rule updates.

\section*{ACKNOWLEDGMENT}

This project was completed as part of the AI Capstone
course at Saint Louis University.

\begin{thebibliography}{00}

\bibitem{b1}
AntLab-Repo,
``Behavioral Model (BMv2),''
[Online]. Available:
\url{https://github.com/p4lang/behavioral-model}.
Accessed: 2026-04-08.

\bibitem{b2}
AntLab-Repo,
``Quark: Implementing Convolutional Neural Networks Entirely on Programmable Data Plane,''
[Online]. Available:
\url{https://github.com/AntLab-Repo/Quark-CNN-P4}.
Accessed: 2026-04-08.

\bibitem{b3}
Moditha Lekkala,
``P4-Based Intrusion Detection System,''
[Online]. Available:
\url{https://github.com/ModithaLekkala/-AI-Capstone-}.
Accessed: 2026-04-08.

\end{thebibliography}

\clearpage

\appendices

\clearpage
\onecolumn

\section*{APPENDIX A\\REPRODUCIBILITY CHECKLIST}

The repository contains the required source code, P4
program, rule file, scripts, documentation, generated plots,
and Docker configuration. The main code is located under
\texttt{src/}, with P4 files in \texttt{src/p4/} and Python testing
scripts in \texttt{src/python/}. The rule file is stored at
\texttt{data/rules.json}. The plots are generated using
\texttt{plots/generate\_plots.py}. The Docker environment is
provided under the \texttt{docker/} directory.

\vspace{0.2cm}

To start the container, run:

\begin{verbatim}
docker compose -f
docker/docker-compose.yml up -d
\end{verbatim}

To run the full benchmark, run:

\begin{verbatim}
cd scripts
powershell -NoProfile
-ExecutionPolicy Bypass
-File .\run_full_benchmark.ps1
\end{verbatim}

The expected benchmark output is:

\begin{verbatim}
ACCURACY: 86/100 tests
correct = 86.0%
TP: 36 TN: 50
FP: 0 FN: 14
Precision: 1.000
Recall : 0.720
\end{verbatim}

To regenerate the plots, run:

\begin{verbatim}
cd plots
python generate_plots.py
\end{verbatim}

The generated plot files are:

\begin{verbatim}
1_roc_curve.png
2_latency_throughput.png
3_mitigation_timeline.png
\end{verbatim}

\clearpage

\section*{APPENDIX B\\TEAM CONTRIBUTION STATEMENT}

This project was completed as a team effort by Venugopal
Rao Ponnamaneni and Moditha Lekkala.

Venugopal Rao Ponnamaneni contributed to the design and
implementation of the P4-based intrusion detection system,
including the data-plane pipeline, rule-based classification
logic, BMv2 switch configuration, testing scripts, traffic
injection process, evaluation metrics, generated plots, final
report writing, and reproducibility documentation.

\vspace{0.2cm}

Moditha Lekkala contributed to project planning, repository
organization, documentation support, initial design discussions,
testing workflow review, validation of system outputs, and
assistance with final report organization. Moditha also
supported the interpretation of benchmark results and helped
refine the project presentation materials.

\vspace{0.2cm}

Both team members collaborated on analyzing the results,
preparing the final capstone submission, and improving the
clarity of the system design and evaluation discussion.

\end{document}